\chapter{ElGamal Encryption}
\label{chap:elGamalEncryption}

We use \gls{eeg} as the semi-homomorphic encrpytion scheme in our voting schemes. 
To understand \gls{eeg} we will first introduce ElGamal Encryption.

\section{ElGamal Encryption}
ElGamal Encryption is an asymmetric encryption scheme. 
Its security is based on the discrete logarithm problem. 
Let $G$ be a cyclic group of order $|n|$ and $g$ be a generator of $G$. 
Then, the discrete logarithm problem states that it is difficult to compute the function $dlog_{(G,g)}:G\to\mathbb{Z}_n, dlog_{(G,g)}(g^{a})=a\:\forall a\in\mathbb{Z}_n$.

Formally, the ElGamal Encryption scheme is defined as follows.
We mostly follow the definition in construction 12.6 of \cite{katz_introduction_2021}.

\begin{definition}
    The ElGamal asymmetric encryption scheme is $\mathcal{S}=(Gen(1^\eta), Enc, Dec)$
\end{definition}


\section{Semi-Homomorphic Encryption Schemes}

\section{Exponential ElGamal Encryption}